% CLAUDE REPORT::WORSTCASE-0001::2026-11-13T10:54Z
% !TEX program = xelatex
% Compiles with XeLaTeX or pdfLaTeX (ASCII source).
\documentclass[11pt]{article}
\usepackage[margin=1in]{geometry}
\usepackage{amsmath,amssymb,amsthm}
\usepackage{booktabs}
\usepackage{xurl}
\usepackage[hidelinks]{hyperref}
\newtheorem{proposition}{Proposition}
\newtheorem{lemma}{Lemma}
\newcommand{\W}{\overline W}
\newcommand{\poly}{\operatorname{poly}}
\title{Worst-case free-strategy derivation cost $\W_R(N)$:\\
quantifier audit of the conditional bridge, exact small Tseitin minima,\\
and the 2025--2026 Res$(\oplus)$ frontier}
\author{Claude scheduled research run (user-directed worst-case directive)}
\date{13 November 2026, 02:54 PST}
\begin{document}
\maketitle

\section*{Status}
\textbf{P versus NP remains OPEN.} No load-bearing arrow is closed in this run.
Statuses used: PROVED/EXACT, KNOWN THEOREM, COMPUTATIONAL, CONDITIONAL,
CANDIDATE, OPEN, FALSIFIED/RETIRED.

\paragraph{Shared target (user directive).}
\[
\W_R(N)=\max_{\substack{F\in\mathrm{UNSAT}\\|\langle F\rangle|\le N}}
\ \min_{\pi:F\vdash_R\bot}\operatorname{cost}_R(\pi),
\qquad
\textbf{(B3)}\quad \mathrm{SAT}\in\mathsf P\ \Longrightarrow\ \W_R(N)\le\poly(N).
\]
Granting the optimal derivation for free is a valid \emph{lower-bound}
relaxation: any procedure that actually outputs an $R$-derivation pays at least
$\min_\pi\operatorname{cost}_R(\pi)$. The load-bearing question is the other
direction: a polynomial-time SAT algorithm is \emph{not} required to output an
$R$-derivation, so B3 is a genuine theorem to be proved, and the corpus
calls it B3 (pointwise and global variants B1/B2).

\section*{CORPUS REVIEWED}
Root \texttt{main.tex} (status line, B1/B2/B3 entries, ER-strength barrier
LSPV-1/592, Haken/Cook PHP controls), \texttt{AGENTS.md}, the November 9
knot/link handoff pointer, \texttt{tasks/} listing, and the newest
\texttt{research/agents/claude/} handoffs (LSPV-1/614, LSPV-2/595,
LSPV-3/610, LSPV-5/523, 01:00--02:24 PST today). Read in full or in the
relevant sections: LSPV-1/334 (free-choice transfer barrier; Tseitin bypass),
LSPV-1/337 (finite reflection calculus $R_T^{\rm len}$), LSPV-1/415
($\Pi^0_1$-completeness of rejection soundness), the codex audits
\texttt{2026-11-11\_1203} and \texttt{2026-11-11\_1331} (quasipolynomial
calibration, clocked $Q_A$, p-optimality, universal reflection
$\mathcal R_{A,\ell}$, Pich--Santhanam witnessing). Repository search for
``relativiz'' located LSPV-1/360 (QBF-oracle B3 nonrelativization) and
LSPV-1/354--367 reflection reports; their conclusions were checked from the
report text, not from summaries.

\section*{DUPLICATION / CONTRADICTION AUDIT}
Before writing, three candidate contributions of this run were found to be
\emph{already in the corpus} and are therefore credited, not re-claimed:
\begin{enumerate}
\item A theory-reflective fixed calculus $R_T$ (proof = algorithm + $T$-proof of
rejection soundness + clocked run) satisfying B3 whenever some polynomial-time
SAT decider is $T$-provably rejection-sound: LSPV-1/337 and LSPV-1/357.
\item The relativization barrier for B3 (a relativizing proof of B3 would apply
to $A^{\mathrm{SAT}}$ and make $\W_R$ unconditionally polynomial): LSPV-1/360.
\item Falsification of pointwise transfer for ordinary Resolution via
Gaussian elimination on odd Tseitin formulas: LSPV-1/334.
\end{enumerate}
No contradiction with root \texttt{main.tex} was found. The phrase ``weakest
sufficient bridge'' in the directive coincides with the corpus' B3.

\section{Quantifier audit of B3 (PROVED/EXACT, elementary)}

\begin{proposition}[Three readings of the bridge]\label{prop:readings}
Let $\mathcal{C}$ range over Cook--Reckhow proof systems.
\begin{enumerate}
\item[(a)] \emph{System chosen after the algorithm:}
$\mathrm{SAT}\in\mathsf P\Rightarrow\exists\mathcal C\,[\W_{\mathcal C}\le\poly]$
is PROVED/EXACT (take $Q_A$: proof of $F$ is $1^{|F|}$, verifier runs $A$).
Its lower-bound partner $\forall\mathcal C\,[\W_{\mathcal C}\ne\poly]$ is
equivalent to $\mathsf{NP}\neq\mathsf{coNP}$ \cite{cookreckhow1979}.
Thus reading (a) routes P$\neq$NP through the strictly stronger-looking
NP$\neq$coNP.
\item[(b)] \emph{System fixed first, lower bound already proved:}
if an unconditional superpolynomial lower bound on $\W_R$ is known
(Resolution, Cutting Planes, Polynomial Calculus, bounded-depth Frege,
tree-like and regular Res$(\oplus)$, \ldots), then B3$(R)$ is logically
equivalent to $\mathrm{SAT}\notin\mathsf P$. B3 for such $R$ is a reformulation
of P$\neq$NP, not a weaker step.
\item[(c)] \emph{System fixed first, no lower bound known:}
B3$(R)$ is a nontrivial statement only here; and it is the case relevant to
the corpus calculus $R_{\rm mix341}$, which p-simulates Extended Resolution
(LSPV-1/592), whose superpolynomial lower bounds are a long-standing open
problem.
\end{enumerate}
\end{proposition}
\begin{proof}
(a) $Q_A$ is sound and complete because $A$ is correct, and verification is
polynomial; Cook--Reckhow: a polynomially bounded system exists iff
NP$=$coNP. (b) B3$(R)\wedge\neg(\W_R\le\poly)$ gives $\mathrm{SAT}\notin\mathsf P$
by contraposition; conversely $\mathrm{SAT}\notin\mathsf P$ makes B3$(R)$
vacuously true. (c) is a statement of where the content lies.
\end{proof}

\begin{lemma}[Squeeze, PROVED/EXACT]\label{lem:squeeze}
For a fixed $R$ the route ``B3$(R)$ $+$ superpolynomial $\W_R$'' requires
\emph{both} that $R$ be strong enough for B3 and weak enough for current lower
bound technology. Every system for which a superpolynomial size lower bound is
known (item (b)) is excluded from proving B3 without simultaneously proving
P$\neq$NP; and the only known positive B3 criterion (provable rejection
soundness, LSPV-1/337,/357) needs $R$ to p-simulate Extended Frege, for which
no superpolynomial lower bound is known.
\end{lemma}
\begin{proof}
First clause: Proposition~\ref{prop:readings}(b). Second clause: Cook's
translation of $\mathrm{PV}$/$S^1_2$ proofs into polynomial-size EF proofs is the
mechanism used by LSPV-1/357; any $R$ admitting it p-simulates EF on the
translated reflection family, and EF superpolynomial lower bounds are OPEN.
\end{proof}
\noindent The squeeze is not an impossibility theorem; it states that
the project's single route currently requires a lower bound at or beyond the
Extended Frege frontier \emph{and} a nonrelativizing B3 (LSPV-1/360).
It sharpens the directive's priority list: item 1 (superpolynomial
$\W_R$) and item 3 (B3) cannot be pursued for different systems.

\section{Survival test: strategy switching is a union system (PROVED/EXACT)}
Free choice of the best sequence already includes adaptive switching
\emph{within} $R$. Switching \emph{between} calculi $R_1,\dots,R_k$ is
modelled by the union system $R_\cup$ (proof $=$ tag $+$ $R_i$-proof), with
$\W_{R_\cup}(N)\le\min_i\W_{R_i}(N)+O(\log k)$ pointwise. Hence a hard family
for the directive's item 4 must be hard simultaneously for Resolution
\emph{and} the GF(2) system (Tseitin dies), \emph{and} Cutting Planes /
Extended Resolution (PHP dies, Cook), \emph{and} algebraic systems. The only
standard candidates conjectured to survive the whole union are random
$k$-CNFs near threshold and certain lifted/combinatorial-design families;
none has a superpolynomial lower bound beyond Res, CP, PC, bounded-depth
Frege, and depth-restricted Res$(\oplus)$ (see Literature).
Status: hardness of any such union for a proposed family is \textbf{CANDIDATE}
at best.

\section{COMPUTATIONAL: exact free-strategy minima on small Tseitin 3-CNFs}
For odd-charge Tseitin formulas on 3-regular graphs (each vertex parity
$\to$ four width-3 clauses), we computed \emph{exactly} the free-strategy
minimum for tree-like Resolution, using the standard equivalence between
tree-like refutations with $L$ leaves and decision trees whose leaves each
falsify an input clause, by dynamic programming over all partial assignments
($f(\rho)=1$ if $\rho$ falsifies a clause, else $\min_x f(\rho_{x=0})+f(\rho_{x=1})$),
and the exact minimum Resolution refutation width by width-bounded saturation.
Each graph is a separate invocation (30-second rule); UNSAT verified by
exhaustive enumeration. Gaussian elimination row-XOR count is the work of
the parity bypass of LSPV-1/334.

\begin{center}
\begin{tabular}{lrrrrrr}
\toprule
Graph & vars & clauses & min tree leaves & min tree size & min width & Gauss XORs\\
\midrule
$K_4$ & 6 & 16 & 28 & 55 & 4 & 5\\
Prism & 9 & 24 & 48 & 95 & 4 & 9\\
$K_{3,3}$ & 9 & 24 & 76 & 151 & 5 & 11\\
$Q_3$ & 12 & 32 & 128 & 255 & n.c. & 20\\
\bottomrule
\end{tabular}
\end{center}
\noindent n.c.: not computed; one-shot saturation exceeded the cap and was not
split this run. Observations (COMPUTATIONAL, tiny $n$, no asymptotic claim):
(i) same vertex/edge counts, different graphs (Prism vs.\ $K_{3,3}$) give
different exact minima, $48$ vs.\ $76$ leaves and width $4$ vs.\ $5$, tracking
expansion as the Ben-Sasson--Wigderson theory predicts; (ii) $Q_3$ needs
$128=2^7$ leaves for 12 variables; (iii) the Gaussian bypass stays at
$\le 20$ row operations. These are calibration data for the free-strategy
$\min_\pi$ in the directive and exact finite optima for a Resolution-family system on Tseitin sources
(LSPV-1/334 explicitly asserted none; this run's search found no earlier
exact tree-like optimum, but the corpus was not exhaustively byte-checked). They are not lower bounds for any stronger
calculus and not SAT-time bounds.
Reproducer: \path{research/agents/claude/worstcase/tseitin_tree_vs_gauss.py},
\path{research/agents/claude/worstcase/tseitin_width.py};
data \path{research/agents/claude/worstcase/results.json}
(SHA-256 \texttt{a20980f1\ldots988c}).

\section{Literature preflight (this run)}
\begin{itemize}
\item \textbf{Optimal proof systems.} Krajíček--Pudlák (1989): NE$=$coNE
$\Rightarrow$ optimal, E$=$NE $\Rightarrow$ p-optimal proof systems; weakened by
Köbler--Messner--Torán. Since SAT$\in\mathsf P\Rightarrow$ E$=$NE, P$=$NP implies
a p-optimal system exists, but \emph{which} system depends on the hypothesis,
so this gives reading (a), not B3 for a fixed $R$. Existence is still OPEN;
Egidy--Glaßer (STOC 2025) and arXiv:2602.02294 (2026) give only oracle
results. KNOWN THEOREM / no transfer to B3.
\item \textbf{Res$(\oplus)$ frontier} (the natural Res$+$XOR union that kills
the Tseitin bypass): no superpolynomial lower bound for unrestricted
dag-like Res$(\oplus)$ is known. 2025 results: Bhattacharya--Chattopadhyay,
exponential bounds for Res$(\oplus)$ of depth nearly $N^2$
(arXiv:2507.23008); Byramji--Impagliazzo, depth $N^{2-\epsilon}$ proofs of an
$O(N)$-clause constant-width CNF need size $\exp(\Omega(N^\epsilon))$
(arXiv:2511.20023); earlier regular Res$(\oplus)$ (Efremenko--Garlík--Itsykson)
and depth $n\log\log n$ (Alekseev--Itsykson). An unverified secondary source
mentions a 2026 depth-$N^{2-\epsilon}$ Tseitin result; not relied on.
\textbf{Directly transferable} only as a ceiling: since $R_{\rm mix341}$
contains Res plus GF(2) row XOR (and DEF), a superpolynomial $\W_{R_{\rm mix341}}$
would in particular settle unrestricted Res$(\oplus)$ on the same family, a
currently open problem.
\item \textbf{Transfer domains.} Coding theory: Tseitin hardness is an
expansion (distance) property; the GF(2) bypass is the decoder exploiting the
same linear structure, an instance of ``hidden algebraic structure is a
weakness'' (AGENTS.md PQ item). Classification: FAILED TRANSFER for any
linear-code-based hard family under a union with GF(2) reasoning.
\end{itemize}

\section*{WHAT CHANGED RELATIVE TO ROOT MAIN}
Nothing load-bearing. Added: (1) explicit three-reading quantifier audit and
squeeze lemma for the directive's B3 (consolidation of LSPV-1/337,/357,/360,
/592 and codex 1331, not a new route); (2) exact tree-like minima and widths
for four Tseitin sources; (3) Res$(\oplus)$ 2025 frontier references.
This run does not edit root \texttt{main.tex}: the scheduled Claude task owns
no \texttt{main.tex} ownership region, so per AGENTS.md the result is left
for LSPV-5 reconciliation.

\section*{ADOPTED / REJECTED / RETIRED PRIOR COMPONENTS}
Adopted: B1/B2/B3 nomenclature; LSPV-1/334 bypass; /337 and /357 reflective
criterion; /360 nonrelativization; /592 ER-strength barrier; codex $Q_A$.
Rejected as a weaker bridge: reading (a) (it is NP$\neq$coNP in disguise).
Retired for B3: any $R$ with a currently proved superpolynomial lower bound
(reading (b)).

\section*{STALE OR OVERSTATED CLAIMS FOUND}
None in the files read. Caution for future writers: the directive's sentence
``finding the optimal sequence can only add work'' is correct only for
procedures that output $R$-derivations; for arbitrary algorithms it presupposes
B3.

\section*{Status table}
\begin{center}
\begin{tabular}{p{0.55\textwidth}p{0.35\textwidth}}
\toprule
Claim & Status\\
\midrule
Proposition~\ref{prop:readings}, Lemma~\ref{lem:squeeze} & PROVED/EXACT (elementary)\\
Union-system bound for strategy switching & PROVED/EXACT\\
Exact tree-like minima / widths, four graphs & COMPUTATIONAL\\
B3 for $R_{\rm mix341}$ or any fixed $R$ with open lower bounds & OPEN\\
Superpolynomial $\W_R$ for any $R\supseteq$ Res$(\oplus)$ & OPEN\\
P versus NP & OPEN\\
\bottomrule
\end{tabular}
\end{center}

\section*{Single next target}
Exact free-strategy minima for the \emph{union} Res$+$GF(2) on a small
non-linear candidate (random 3-CNF at ratio $\approx 4.26$, $n\le 12$),
tree-like first via the same DP with XOR queries added, to see whether any
candidate shows growth that the GF(2) bypass does not erase.

\begin{thebibliography}{9}
\bibitem{cookreckhow1979} S.~A. Cook, R.~A. Reckhow, The relative efficiency of
propositional proof systems, \emph{J.\ Symb.\ Logic} 44 (1979) 36--50.
\bibitem{kp1989} J.~Krajíček, P.~Pudlák, Propositional proof systems, the
consistency of first order theories and the complexity of computations,
\emph{J.\ Symb.\ Logic} 54 (1989) 1063--1079.
\bibitem{bsw2001} E.~Ben-Sasson, A.~Wigderson, Short proofs are narrow ---
resolution made simple, \emph{J.\ ACM} 48 (2001) 149--169.
\bibitem{urquhart1987} A.~Urquhart, Hard examples for resolution,
\emph{J.\ ACM} 34 (1987) 209--219.
\bibitem{bc2025} S.~Bhattacharya, A.~Chattopadhyay, Res$(\oplus)$ lower bounds
for depth nearly $N^2$, arXiv:2507.23008 (2025).
\bibitem{bi2025} S.~Byramji, R.~Impagliazzo, arXiv:2511.20023 (2025).
\bibitem{eg2025} F.~Egidy, C.~Glaßer, oracle constructions for optimal proof
systems, STOC 2025; see also arXiv:2602.02294 (2026).
\end{thebibliography}
\end{document}
% END_OF_VERIFIED_REPORT::CLAUDE::WORSTCASE-0001
